% Part: second-order-logic
% Chapter: syntax-and-semantics
% Section: terms-formulas

\documentclass[../../../include/open-logic-section]{subfiles}

\begin{document}

% SEGMENT OLP-0325-S01

\olfileid{sol}{syn}{frm}

\olsection{உறுப்புகளும் \printtoken{P}{formula}}

முதல்தரத் தருக்கத்தைப் போலவே, இரண்டாம்தரத் தருக்கத்தின் கோவைகள்
\emph{!!{variable}s}, \emph{!!{constant}s}, \emph{!!{predicate}s},
சில நேரங்களில் \emph{!!{function}s} ஆகியவற்றைக் கொண்ட அடிப்படைச்
சொற்கோவையிலிருந்து கட்டியெழுப்பப்படுகின்றன. இவற்றுடன் தருக்க
இணைப்பிகள், அளவையடைகள், அடைப்புக்குறிகள் மற்றும் காற்புள்ளிகள் போன்ற
நிறுத்தற்குறிகளைச் சேர்த்து, \emph{உறுப்புகளும்}
\emph{!!{formula}s உம்} அமைக்கப்படுகின்றன. பொருள்களுக்கான மாறிகளோடு,
உறவுகளுக்கும் சார்புகளுக்குமான மாறிகளையும் இரண்டாம்தரத் தருக்கம்
கொண்டிருப்பதும், அவற்றின்மீது அளவையிடலை அனுமதிப்பதுமே வேறுபாடு.
எனவே இரண்டாம்தரத் தருக்கத்தின் தருக்கக் குறிகள் முதல்தரத்
தருக்கத்தின் குறிகளுடன் பின்வருவனவற்றையும் கொண்டவை:

\begin{enumerate}
\item ஒவ்வோர் இட எண்ணிக்கை~$n$ க்கும் இரண்டாம்தர உறவு
  !!{variable}s இன் ஓர் !!{denumerable}s கணம்:
  $\Obj V_0^n$, $\Obj V_1^n$, $\Obj V_2^n$, \dots
\item இரண்டாம்தரச் சார்பு !!{variable}s இன் ஓர் !!{denumerable}s கணம்:
  $\Obj u_0^n$, $\Obj u_1^n$, $\Obj u_2^n$, \dots
\end{enumerate}

முதல்தர மாறிகள் $\Obj v_i$ க்கான மேல்மொழி மாறிகளாக $x$, $y$, $z$
ஆகியவற்றைப் பயன்படுத்துவது போல, $\Obj V_i^n$ க்கான மேல்மொழி
மாறிகளாக $X$, $Y$, $Z$ முதலியவற்றையும்,~$\Obj u_i^n$ க்கான
மேல்மொழி மாறிகளாக $u$, $v$ முதலியவற்றையும் பயன்படுத்துவோம்.

\begin{explain}
ஒரு முதல்தர மொழியைக் குறிப்பிடும் அதே முறையில், அதன் !!{constant}s,
!!{function}s, !!{predicate}s ஆகியவற்றைப் பட்டியலிட்டு, இரண்டாம்தர
மொழியின் தருக்கமல்லாத குறிகளும் குறிப்பிடப்படுகின்றன.

முதல்தரத் தருக்கத்தில் !!{identity}~$\eq$ பொதுவாகச் சேர்க்கப்படுகிறது.
முதல்தரத் தருக்கத்தில், ஒரு மொழி~$\Lang{L}$ இன் தருக்கமல்லாத குறிகள்
சுவையான எதையும் வெளிப்படுத்த இன்றியமையாதவை. தருக்கமல்லாத குறிகள்
எதையும் பயன்படுத்தாத !!{sentence}s நிச்சயமாக உள்ளன; ஆனால்~$\eq$
மட்டும் கொண்டு சுவையான எதையும் கூறுவது கடினம். இரண்டாம்தரத்
தருக்கத்தில், உறவு மற்றும் சார்பு மாறிகள் வரம்பின்றிக் கிடைப்பதால்,
தருக்கமல்லாத குறிகளின் தனிப்பட்ட தொகுதி இல்லாமலேயே ஒரு முதல்தர
மொழியில் சொல்லக்கூடிய எதையும் நம்மால் சொல்ல முடியும்.
\end{explain}

\begin{defn}[இரண்டாம்தர உறுப்புகள்]
~$\Lang L$ இன் \emph{இரண்டாம்தர உறுப்புகளின்} கணமான $\TrmSOL[L]$,
\olref[fol][syn][frm]{defn:terms} உடன் பின்வரும் விதியைச் சேர்ப்பதன்
மூலம் வரையறுக்கப்படுகிறது:
\begin{enumerate}
\item $u$ ஓர் $n$-இடச் சார்பு மாறியாகவும், $t_1$, \dots, $t_n$
  உறுப்புகளாகவும் இருந்தால், $\Atom{u}{t_1, \ldots, t_n}$ ஓர் உறுப்பு.
\end{enumerate}
\end{defn}

\begin{explain}
எனவே, ஓர் இரண்டாம்தர உறுப்பு முதல்தர உறுப்பைப் போலவே தோன்றும்.
ஆனால், முதல்தர உறுப்பில் !!a{function}~$\Obj{f^n_i}$ வரும் இடத்தில்,
இரண்டாம்தர உறுப்பில் சார்பு மாறி~$\Obj{u^n_i}$ வரலாம்.
\end{explain}

\begin{defn}[இரண்டாம்தர \usetoken{s}{formula}]
மொழி~$\Lang L$ இன் \emph{இரண்டாம்தர !!{formula}s} கணமான
~$\FrmSOL[L]$, \olref[fol][syn][frm]{defn:formulas} உடன் பின்வரும்
விதிகளைச் சேர்ப்பதன் மூலம் வரையறுக்கப்படுகிறது:
\begin{enumerate}
\item $X$ ஓர் $n$-இடப் பயனிலை மாறியாகவும், $t_1$, \dots, $t_n$
  ஆகியவை~$\Lang L$ இன் இரண்டாம்தர உறுப்புகளாகவும் இருந்தால்,
  $\Atom{X}{t_1,\ldots, t_n}$ ஓர் அணு !!{formula}.

\tagitem{prvAll}{$!A$ !!a{formula} ஆகவும் $u$ ஒரு சார்பு மாறியாகவும்
  இருந்தால், $\lforall[u][!A]$ !!a{formula}.}{}

\tagitem{prvAll}{$!A$ !!a{formula} ஆகவும் $X$ ஒரு பயனிலை மாறியாகவும்
  இருந்தால், $\lforall[X][!A]$ !!a{formula}.}{}

\tagitem{prvEx}{$!A$ !!a{formula} ஆகவும் $u$ ஒரு சார்பு மாறியாகவும்
  இருந்தால், $\lexists[u][!A]$ !!a{formula}.}{}

\tagitem{prvEx}{$!A$ !!a{formula} ஆகவும் $X$ ஒரு பயனிலை மாறியாகவும்
  இருந்தால், $\lexists[X][!A]$ !!a{formula}.}{}
\end{enumerate}
\end{defn}

\end{document}
